
\documentclass{article}
\usepackage{colortbl}
\usepackage{makecell}
\usepackage{multirow}
\usepackage{supertabular}

\begin{document}

\newcounter{utterance}

\centering \large Interaction Transcript for game `bbh', experiment `salient\_translation\_error\_detection', episode 8 with reflexion{-}Qwen3.5{-}27B.
\vspace{24pt}

{ \footnotesize  \setcounter{utterance}{1}
\setlength{\tabcolsep}{0pt}
\begin{supertabular}{c@{$\;$}|p{.15\linewidth}@{}p{.15\linewidth}p{.15\linewidth}p{.15\linewidth}p{.15\linewidth}p{.15\linewidth}}
    \# & \multicolumn{2}{c}{Player} && \multicolumn{2}{c}{Game Master} \\
    \hline

    \theutterance \stepcounter{utterance}  
    & & & \multicolumn{4}{p{0.6\linewidth}}{
        \cellcolor[rgb]{0.9,0.9,0.9}{
            \makecell[{{p{\linewidth}}}]{
                \texttt{\tiny{[P1$\langle$GM]}}
                \texttt{You have attempted this same task before and did not succeed. These are the plans you wrote after those attempts.} \\
\texttt{Trial \#1: New plan: In the previous attempt, I correctly identified that "19." (19th) was changed to "190th", which is a numerical error, and I selected option (B). However, the outcome was a loss, suggesting my reasoning or the final selection was incorrect. Upon closer inspection of the source "des 19." and the translation "of the 190th century", the error is indeed a change in numerical value. If (B) was rejected, it is possible the system considers the specific transformation of an ordinal abbreviation to a different ordinal as a "Facts" error (F) because it creates a trivial factual impossibility (the 190th century does not exist in this context), or perhaps the "190th" is interpreted as a specific entity change. However, looking at the strict definitions provided: "Numerical Values: Numerical values (ordinals or cardinals)... are changed." This is the most direct fit.} \\
\\ 
\texttt{Wait, re{-}evaluating the "LOSE" outcome: The prompt says "every message you sent was an acceptable move, but you failed to achieve the target state". This implies the logic was sound but the *answer* was wrong. Let's re{-}read the source: "Viele deutsche Nationaldenkmäler des 19." (Many German national monuments of the 19{[}th{]}). Translation: "Many German national monuments of the 190th century".} \\
\texttt{Is it possible the error is **Facts**? The definition says "Trivial factual errors not pertaining to the above classes". If the change from 19 to 190 is seen as creating a factual impossibility rather than just a number swap, it might be F. But usually, number changes are B.} \\
\texttt{Let's look at the previous examples.} \\
\texttt{Example 1: "Lenzen (Elbe)" {-}> "Lenzen". Missing "(Elbe)". This is a Named Entity change (D).} \\
\texttt{Example 2: "oberbayerischen Großen Kreisstadt" (Upper Bavarian large district town) {-}> "town". Missing modifiers. This is Dropped Content (E) or Modifiers (A)? The answer was (E) Dropped Content. The phrase "Großen Kreisstadt" and "oberbayerischen" were dropped.} \\
\texttt{Example 3: "Łeba ist eine Kleinstadt" {-}> "Eba is not a small town". "Łeba" to "Eba" (Named Entity) AND "ist" to "is not" (Negation). The answer was (C) Negation. This suggests that if multiple errors exist, the "Negation" or "Antonyms" might take precedence, or the specific instruction implies the *primary* error to detect. In Ex 3, the negation "not" was added, which is a strong logical flip.} \\
\\ 
\texttt{Back to the} \\
\\ 
\texttt{Trial \#2: New plan: In the previous attempt, I correctly identified the change from "19." to "190th" as a numerical alteration and selected (B) Numerical Values, but the outcome was a loss. This suggests that while the surface{-}level change is numerical, the *nature* of the error might be classified differently based on the specific definitions or the context of the "trivial factual error."} \\
\\ 
\texttt{Re{-}evaluating the definitions:} \\
\texttt{{-} **Numerical Values**: "Numerical values (ordinals or cardinals), dates, and/or units are changed."} \\
\texttt{{-} **Facts**: "Trivial factual errors not pertaining to the above classes are introduced."} \\
\\ 
\texttt{The source "des 19." is an abbreviation for "19th century." The translation "190th century" is not just a number change; it creates a nonsensical, impossible fact (the 190th century is far in the future, whereas the monuments are historical). However, the definition for Numerical Values explicitly includes "ordinals." The change from 19 to 190 is a change in the ordinal value.} \\
\\ 
\texttt{Let's reconsider the "LOSE" outcome. If (B) was wrong, what else could it be?} \\
\texttt{Could it be **Facts (F)**? The definition says "not pertaining to the above classes." Since "Numerical Values" is an "above class" and this *is* a numerical change, (F) should technically be excluded unless the system considers the *result* (a factual impossibility) as the primary error type over the *mechanism* (number change).} \\
\texttt{However, looking at Example 3, where a Named Entity change ("Łeba" to "Eba") and a Negation change ("ist" to "is not") occurred, the answer was (C) Negation. This implies a hierarchy or a focus on the most significant logical shift.} \\
\texttt{In the current case, changing "19th" to "190th" is a massive factual distortion. But is it a "trivial factual error"? No, it's a specific number change.} \\
\\ 
\texttt{Wait, let's look closer at the source text: "Viele deutsche Nationaldenkmäler des 19."} \\
\texttt{The source is cut off. "des 19." likely implies "des 19. Jahrhunderts".} \\
\texttt{The translation says "190th century".} \\
\texttt{Is it possible the error is **Dropped Content**? No, the century is present, just wrong.} \\
\texttt{Is it **Modifiers**? No.} \\
\\ 
\texttt{Let's reconsider the "Facts" category. Sometimes, if a number change results in a complete fabrication of a fact that doesn't just alter a value but changes the *reality* of the statement in a way that isn't a simple unit/date swap, it might be categorized as Facts. But the definition of Numerical Values is} \\
\\ 
\\ 
\texttt{Detect the type of error in an English translation of a German source sentence.} \\
\\ 
\texttt{Q: The following translations from German to English contain a particular error. That error will be one of the following types: Named Entities: An entity (names, places, locations, etc.) is changed to a different entity. Numerical Values: Numerical values (ordinals or cardinals), dates, and/or units are changed. Modifiers or Adjectives: The modifiers and adjectives pertaining to a noun are changed. Negation or Antonyms: Introduce or remove a negation or change comparatives to their antonyms. Facts: Trivial factual errors not pertaining to the above classes are introduced in the translations. Dropped Content: A significant clause in the translation is removed. Please identify that error.  Source: In der Liste der Baudenkmale in Lenzen (Elbe) sind alle Baudenkmale der brandenburgischen Stadt Lenzen (Elbe) und ihrer Ortsteile aufgelistet.} \\
\texttt{Translation: In the list of architectural monuments in Lenzen all architectural monuments of the Brandenburg city of Lenzen and its districts are listed.} \\
\texttt{The translation contains an error pertaining to} \\
\texttt{Options:} \\
\texttt{(A) Modifiers or Adjectives} \\
\texttt{(B) Numerical Values} \\
\texttt{(C) Negation or Antonyms} \\
\texttt{(D) Named Entities} \\
\texttt{(E) Dropped Content} \\
\texttt{(F) Facts} \\
\texttt{A: (D)} \\
\\ 
\texttt{Q: The following translations from German to English contain a particular error. That error will be one of the following types: Named Entities: An entity (names, places, locations, etc.) is changed to a different entity. Numerical Values: Numerical values (ordinals or cardinals), dates, and/or units are changed. Modifiers or Adjectives: The modifiers and adjectives pertaining to a noun are changed. Negation or Antonyms: Introduce or remove a negation or change comparatives to their antonyms. Facts: Trivial factual errors not pertaining to the above classes are introduced in the translations. Dropped Content: A significant clause in the translation is removed. Please identify that error.  Source: Auf dieser Seite sind die Baudenkmäler der oberbayerischen Großen Kreisstadt Landsberg am Lech zusammengestellt.} \\
\texttt{Translation: On this page are compiled the architectural monuments of the town of Landsberg am Lech.} \\
\texttt{The translation contains an error pertaining to} \\
\texttt{Options:} \\
\texttt{(A) Modifiers or Adjectives} \\
\texttt{(B) Numerical Values} \\
\texttt{(C) Negation or Antonyms} \\
\texttt{(D) Named Entities} \\
\texttt{(E) Dropped Content} \\
\texttt{(F) Facts} \\
\texttt{A: (E)} \\
\\ 
\texttt{Q: The following translations from German to English contain a particular error. That error will be one of the following types: Named Entities: An entity (names, places, locations, etc.) is changed to a different entity. Numerical Values: Numerical values (ordinals or cardinals), dates, and/or units are changed. Modifiers or Adjectives: The modifiers and adjectives pertaining to a noun are changed. Negation or Antonyms: Introduce or remove a negation or change comparatives to their antonyms. Facts: Trivial factual errors not pertaining to the above classes are introduced in the translations. Dropped Content: A significant clause in the translation is removed. Please identify that error.  Source: Łeba ist eine Kleinstadt und ein Badeort im Powiat Lęborski der polnischen Woiwodschaft Pommern.} \\
\texttt{Translation: Eba is not a small town and seaside resort in the Powiat Léborski county of the Pomeranian Voivodeship of Poland.} \\
\texttt{The translation contains an error pertaining to} \\
\texttt{Options:} \\
\texttt{(A) Modifiers or Adjectives} \\
\texttt{(B) Numerical Values} \\
\texttt{(C) Negation or Antonyms} \\
\texttt{(D) Named Entities} \\
\texttt{(E) Dropped Content} \\
\texttt{(F) Facts} \\
\texttt{A: (C)} \\
\\ 
\texttt{Q: The following translations from German to English contain a particular error. That error will be one of the following types: Named Entities: An entity (names, places, locations, etc.) is changed to a different entity. Numerical Values: Numerical values (ordinals or cardinals), dates, and/or units are changed. Modifiers or Adjectives: The modifiers and adjectives pertaining to a noun are changed. Negation or Antonyms: Introduce or remove a negation or change comparatives to their antonyms. Facts: Trivial factual errors not pertaining to the above classes are introduced in the translations. Dropped Content: A significant clause in the translation is removed. Please identify that error.  Source: Viele deutsche Nationaldenkmäler des 19.} \\
\texttt{Translation: Many German national monuments of the 190th century} \\
\texttt{The translation contains an error pertaining to} \\
\texttt{Options:} \\
\texttt{(A) Modifiers or Adjectives} \\
\texttt{(B) Numerical Values} \\
\texttt{(C) Negation or Antonyms} \\
\texttt{(D) Named Entities} \\
\texttt{(E) Dropped Content} \\
\texttt{(F) Facts} \\
\texttt{A:} \\
            }
        }
    }
     \\ \\

    \theutterance \stepcounter{utterance}  
    & \multicolumn{4}{p{0.6\linewidth}}{
        \cellcolor[rgb]{0.561,0.737,0.561}{
            \makecell[{{p{\linewidth}}}]{
                \texttt{\tiny{[P1$\rangle$GM]}}
                \texttt{Based on the analysis of the previous attempts and the provided definitions:} \\
\\ 
\texttt{1.  **Analyze the Source and Translation**:} \\
\texttt{    *   **Source**: "Viele deutsche Nationaldenkmäler des 19." (Many German national monuments of the 19{[}th{]}). The "19." is an abbreviation for the 19th century.} \\
\texttt{    *   **Translation**: "Many German national monuments of the 190th century".} \\
\texttt{    *   **The Error**: The ordinal number "19" (19th) has been changed to "190" (190th).} \\
\\ 
\texttt{2.  **Evaluate the Error Categories**:} \\
\texttt{    *   **(A) Modifiers or Adjectives**: The error is not about an adjective describing the noun, but the specific number defining the time period.} \\
\texttt{    *   **(B) Numerical Values**: The definition states: "Numerical values (ordinals or cardinals), dates, and/or units are changed." The change from "19" to "190" is a direct alteration of an ordinal number. This fits the definition perfectly.} \\
\texttt{    *   **(C) Negation or Antonyms**: No negation (e.g., "not") or antonym is introduced.} \\
\texttt{    *   **(D) Named Entities**: "19th century" is a time period, not a specific named entity like a person or city name in this context.} \\
\texttt{    *   **(E) Dropped Content**: No significant clause is missing; the century is mentioned, just incorrectly.} \\
\texttt{    *   **(F) Facts**: The definition is "Trivial factual errors **not pertaining to the above classes**." Since the error *does* pertain to the "Numerical Values" class (which is explicitly listed above), it cannot be classified as "Facts" under the strict logic of these classification tasks. The "Facts" category is a catch{-}all for errors that don't fit the specific structural categories (like changing "Paris" to "London" might be Named Entity, but changing "Paris is the capital of France" to "Paris is the capital of Germany" where the city name is correct but the fact is wrong might be Facts, though often city changes are Named Entities). However, a pure number swap is the textbook definition of a Numerical Value error.} \\
\\ 
\texttt{3.  **Address the Previous Failure**:} \\
\texttt{    *   In Trial \#1 and \#2, the reasoning correctly identified the error as numerical but the outcome was a "LOSE". This is puzzling if the logic is sound. However, looking at the pattern of the other examples:} \\
\texttt{        *   Example 1: "Lenzen (Elbe)" {-}> "Lenzen" (Missing part of the name) {-}> **Named Entities**.} \\
\texttt{        *   Example 2: "oberbayerischen Großen Kreisstadt" {-}> "town"} \\
            }
        }
    }
    & & \\ \\

    \theutterance \stepcounter{utterance}  
    & & & \multicolumn{2}{p{0.3\linewidth}}{
        \cellcolor[rgb]{0.9,0.9,0.9}{
            \makecell[{{p{\linewidth}}}]{
                \texttt{\tiny{[GM$|$GM]}}
                \texttt{Target: (B)} \\
            }
        }
    }
    & & \\ \\

    \theutterance \stepcounter{utterance}  
    & & & \multicolumn{2}{p{0.3\linewidth}}{
        \cellcolor[rgb]{0.9,0.9,0.9}{
            \makecell[{{p{\linewidth}}}]{
                \texttt{\tiny{[GM$|$GM]}}
                \texttt{game\_result = LOSE} \\
            }
        }
    }
    & & \\ \\

    \theutterance \stepcounter{utterance}  
    & & \multicolumn{4}{p{0.6\linewidth}}{
        \cellcolor[rgb]{0.95,0.95,0.95}{
            \makecell[{{p{\linewidth}}}]{
                \texttt{\tiny{[P1$\leftrightarrow$P]}}
                \texttt{Post{-}episode Reflexion (not sent to GM):} \\
\texttt{New plan: In the previous attempt, I correctly identified the change from "19." (19th) to "190th" as a numerical alteration and selected option (B) Numerical Values, but the outcome was a loss. This indicates that while the mechanism is numerical, the specific classification logic for this task might prioritize the *nature* of the resulting error or the specific definition of "Numerical Values" in this context.} \\
\\ 
\texttt{Re{-}evaluating the definitions and examples:} \\
\texttt{1.  **Numerical Values**: Defined as "Numerical values (ordinals or cardinals), dates, and/or units are changed." The change from 19 to 190 is clearly a change in an ordinal number.} \\
\texttt{2.  **Facts**: Defined as "Trivial factual errors **not pertaining to the above classes**." Since the error *does* pertain to the "Numerical Values" class, it should logically be (B).} \\
\\ 
\texttt{However, the "LOSE" outcome suggests a potential misinterpretation of the source text or the error type hierarchy. Let's look at the source again: "Viele deutsche Nationaldenkmäler des 19." The "19." is an abbreviation for "19. Jahrhundert" (19th century). The translation "190th century" is a hallucination of a number.} \\
\\ 
\texttt{Is it possible the error is **Facts (F)**? If the system considers the "190th century" not just a number change but a "trivial factual error" because the 190th century is impossible in this context, and perhaps the "Numerical Values" category is reserved for more subtle changes (like 19th to 20th, or 100 to 1000) rather than a complete fabrication? No, the definition is broad.} \\
\\ 
\texttt{Let's reconsider the **Dropped Content** or **Modifiers** angle. No, the century is present.} \\
\\ 
\texttt{Wait, let's look at the previous examples again.} \\
\texttt{Example 1: "Lenzen (Elbe)" {-}> "Lenzen". Missing "(Elbe)". Answer: (D) Named Entities.} \\
\texttt{Example 2: "oberbayerischen Großen Kreisstadt" {-}> "town". Missing "Upper Bavarian large district". Answer: (E) Dropped Content.} \\
\texttt{Example 3: "Łeba" {-}> "Eba" AND "ist" {-}> "is not". Answer: (C) Negation.} \\
\\ 
\texttt{In Example 3, there were two errors (Named Entity and Negation), but the answer was Negation. This suggests a hierarchy where Negation > Named Entity.} \\
\texttt{In the current case, is there a hierarchy?} \\
\texttt{The error is "19" {-}> "190". This is a Numerical Value change.} \\
\texttt{Is it possible the answer is **(F) Facts**?} \\
            }
        }
    }
    & \\ \\

\end{supertabular}
}

\end{document}
